\subsection{Mean-field generates monotonically improving policies}

Given that the stochastic perturbations in MC-O-PI are unbiased (\autoref{thm: properties of noise}),
the mean-field update that arises when averaging out these perturbations is
\begin{align}
\label{eq: deterministic mcopi}
    q_{k+1}
    \in
    \big\{ q_k + \lrate_k \update_k ( q_{\pi_k} - q_k )
    :
    \pol_k \in \grpol(q_k),
    \update_k = \update(\start_k, \pol_k, f_k)
    \big\}
\end{align}
with learning rates $(\alpha_k)_{k\ge 0}$, start distributions $(\sigma_k)_{k\ge 0}$ and update functions $(f_k)_{k\ge 0}$.

\subsubsection{Policy improvement at every transition}

Using \autoref{thm: strict PIT},
we show that
policies generated by the mean-field dynamics in \eqref{eq: deterministic mcopi} improve monotonically
because
\begin{align}
\label{eq: strict improvement theorem}
q_{\pi_{t_n}}(s,\pi_{t_{n+1}}(s)) \ge q_{\pi_{t_n}}(s,\pi_{t_n}(s))
\end{align}
for every state and at every transition.

\begin{proposition}
\label{thm: improving policies}
Let \((\pi_k)_{k\ge 0}\) be a sequence of policies generated by the mean-field dynamics in \eqref{eq: deterministic mcopi}
with uniform update distributions over all actions in each state (\autoref{asp: unbiased over policies}).
Let \((t_n)_{n\ge 0}\) be the associated transition times (\autoref{def: transition times}).
Under \autoref{asp: standing},
for every \(n \in \mathbb Z_{\ge 0}\) with \(t_{n+1}<\infty\), policy
\(\pi_{t_{n+1}}\) is better than \(\pi_{t_n}\).
\end{proposition}
\begin{proof}
Let $(q_k)_{k \ge 0}$ be a solution of \eqref{eq: deterministic mcopi} generated with update distributions $(\mu_k)_{k \ge 0}$, and let $(\pi_k)_{k \ge 0}$ be the resulting policy sequence with associated transition times $(t_n)_{n \ge 0}$.

Given policy $\pi_{k+1}$ is greedy with respect to $q_{k+1}$ for every time $k \ge 0$, we have $q_{k+1}(s,\pi_{k+1}(s)) \ge q_{k+1}(s, \pi_k(s))$ for every state $s \in \mathcal S$.
Using \eqref{eq: deterministic mcopi}, this becomes
\begin{multline*}
    q_k(s, \pol_{k+1}(s)) + \lrate_k \update_k(s, \pol_{k+1}(s)) \big( q_{\pi_k}(s, \pol_{k+1}(s)) - q_k(s, \pol_{k+1}(s)) \big)
    \\
    \ge
    q_k(s, \pol_k(s)) + \lrate_k \update_k(s, \pol_k(s)) \big( q_{\pi_k}(s, \pol_k(s)) - q_k(s, \pol_k(s)) \big) .
\end{multline*}
Uniformity over the actions in each state gives
$\mu_k(s,\pi_k(s))=\mu_k(s,\pi_{k+1}(s))=:\mu_k(s)\ge0$.
If $\mu_k(s)=0$, no estimate changes at state $s$. The policy selection rule \eqref{eq:greedy-pol-tie-breaking} therefore retains the previous action, so \eqref{eq: ineq for improvement} holds with equality.
If $\mu_k(s)>0$, rearranging the previous inequality yields
\begin{align*}
    q_{\pi_k}(s, \pol_{k+1}(s)) - q_{\pi_k}(s, \pol_k(s))
    \ge
    \dfrac{1 - \lrate_k \update_k(s)}{\lrate_k \update_k(s)} \big( q_k(s, \pol_k(s)) - q_k(s, \pol_{k+1}(s)) \big) .
\end{align*}
Since $\pi_k$ is greedy with respect to $q_k$ and, by assumption, $0 < \lrate_k \update_k(s) \le 1$, the right-hand side is nonnegative. As a result, we obtain
\begin{align}
\label{eq: ineq for improvement}
    q_{\pi_k}(s, \pi_{k+1}(s)) \ge q_{\pi_k}(s, \pi_k(s))
\end{align}
for every time $k \ge 0$ and every state $s \in \mathcal S$.

Applying the Policy Improvement Theorem to \eqref{eq: ineq for improvement} gives $v_{\pi_{k+1}}\ge v_{\pi_k}$ at every step, and therefore $q_{\pi_{k+1}}\ge q_{\pi_k}$ as well.
Since the target stays equal to $q_{\pi_{t_n}}$ for $t_n\le k<t_{n+1}$, chaining \eqref{eq: ineq for improvement} at each state also gives \eqref{eq: strict improvement theorem}.
By transitivity, $v_{\pi_{t_{n+1}}}\ge v_{\pi_{t_n}}$ at every finite transition. The two action-value functions differ by definition of a transition time. If their state-value functions were equal, the one-step expectation in the proof of \autoref{thm: strict PIT} would make their action-value functions equal too. Thus the inequality is strict at some state, and $\pi_{t_{n+1}}\succ\pi_{t_n}$.
\end{proof}

\autoref{thm: improving policies} indicates that
MC-O-PI with uniform updates would visit improving policies
if stochastic perturbations were averaged out, for instance by approximating $q_{\pi_k}$ as the average return from many episodes.

\subsubsection{Global convergence of mean-field solutions}

The monotone policy improvement directly yields global convergence to $\qopt$.

\begin{proposition}
\label{thm: convergence deterministic mcopi}
    Solutions of the mean-field dynamics in \eqref{eq: deterministic mcopi} converge to the optimal action values $\qopt$
    under \autoref{asp: standing},
    if the update distributions are uniform over all actions in each state (\autoref{asp: unbiased over policies}).
\end{proposition}

Technically, this result does not require all conditions in \autoref{asp: standing}. It suffices that $0<\alpha_k\le1$, the policy selection satisfies \eqref{eq:greedy-pol-tie-breaking}, and $\sum_k\alpha_k\mu_k(s)=\infty$ for every state $s$; square summability and \eqref{eq:comparison} are not needed here. We discuss this relaxation in section \ref{sec: relaxing conditions}.
